% ======================================================================
\section{Comparative assessment}
\label{sec:assessment}
% ======================================================================

\subsection{Summary of properties}

\Cref{tab:summary} summarizes the families reviewed above against the
properties (P1)--(P7) of \cref{sec:desiderata}. The entries are qualitative and
reflect typical behaviour reported in the literature and in our computed
examples; any given implementation can do better or worse.

\begin{table}[tbp]
\centering
\caption{Qualitative comparison of dimensionality reduction families for
fluid-dynamics machine learning. Columns: \emph{Dec.}\ explicit, cheap decoder;
\emph{Transp.}\ latent size needed for transport-dominated or multi-scale flows
(slowly decaying Kolmogorov $n$-width); \emph{Data} number of snapshots needed
to build the representation; \emph{Scal.}\ cost and memory for large $N$ and
mesh flexibility; \emph{Interp.}\ interpretability; \emph{Acc.}\ attainable
accuracy as the latent size grows. $+$ favourable, $\circ$ intermediate,
$-$ unfavourable.}
\label{tab:summary}
\footnotesize
\setlength{\tabcolsep}{3pt}
\renewcommand{\arraystretch}{1.15}
\begin{tabularx}{\textwidth}{@{}>{\raggedright\arraybackslash}p{4.3cm} *{6}{c} >{\raggedright\arraybackslash}X@{}}
\toprule
Family & Dec. & Transp. & Data & Scal. & Interp. & Acc. & Representative fluid uses \\
\midrule
POD / PCA & $+$ & $-$ & $\circ$ & $+$ & $+$ & $+$ & Galerkin ROMs, POD--LSTM/transformer, OpInf \\
DMD, SPOD, balanced POD & $+$ & $-$ & $\circ$ & $+$ & $+$ & $+$ & Modal analysis, Koopman models, control \\
Transport-adapted linear & $+$ & $\circ$ & $\circ$ & $\circ$ & $+$ & $+$ & Shocks, fronts, traveling waves \\
Sparse, clustered, local ROBs & $\circ$ & $\circ$ & $\circ$ & $\circ$ & $+$ & $+$ & Multi-scale Galerkin ROMs, regime switching \\
Domain decomposition, GFEM, DLS & $+$ & $\circ$ & $+$ & $+$ & $+$ & $+$ & Extreme-scale ROMs, turbulent latent spaces \\
Kernel / graph manifold learning & $-$ & $+$ & $\circ$ & $-$ & $\circ$ & $\circ$ & Manifold identification, regime analysis \\
Quadratic manifolds & $+$ & $\circ$ & $\circ$ & $+$ & $\circ$ & $+$ & Projection ROMs beyond linear subspaces \\
MLP / convolutional AE & $+$ & $+$ & $-$ & $\circ$ & $-$ & $-$ & Nonlinear ROMs, AE--LSTM/transformer \\
VAE, $\beta$-VAE, VQ-VAE & $+$ & $+$ & $-$ & $\circ$ & $\circ$ & $-$ & Disentangled modes, compression, generation \\
Graph autoencoders & $+$ & $+$ & $-$ & $\circ$ & $\circ$ & $-$ & Unstructured meshes \\
Implicit neural representations & $+$ & $+$ & $-$ & $\circ$ & $-$ & $-$ & Mesh-agnostic ROMs, generative turbulence \\
Latent neural operators & $+$ & $+$ & $-$ & $\circ$ & $-$ & $-$ & Parametric surrogates, general geometries \\
\bottomrule
\end{tabularx}
\end{table}

Two broad patterns emerge. First, representations with linear, explicit
encoders and decoders (POD, DMD, GFEM/DLS) are cheap, reproducible, and reach
arbitrarily high accuracy as the latent size grows. They are the natural choice
when training data are limited or when accuracy of small scales matters.
Learned nonlinear decoders achieve much smaller latent spaces for
transport-dominated and low-intrinsic-dimension problems but saturate in
accuracy and require more data and tuning. Second, the right choice depends on
the \emph{flow}, not only on the method. \Cref{fig:ldc_pod} shows that the
$\mathrm{Re}=15{,}000$ cavity is dominated by a few coherent modes, so global
POD is close to optimal for reconstructing it at a fixed latent size
(\cref{fig:dls_compare}a). The same is not true for the advected pulse of
\cref{fig:nwidth,fig:ae_vs_pod} or, by extension, for strongly convective or
broadband turbulent flows.

\subsection{Evidence from comparative studies}

Systematic comparisons remain scarce, and their conclusions are nuanced.
\textcite{csala2022comparing} and \textcite{mendez2023linear} found that
nonlinear methods can represent canonical flows with fewer coordinates than
PCA, but that the advantage depends on the flow, the arrangement of the data,
and the task (filtering, pattern identification, or compression).
\textcite{gruber2022comparison} showed that no single autoencoder architecture
dominates: the best choice depends on latent size and data structure.
\textcite{swischuk2019projection} compared several regression methods on POD
coefficients and argued for retaining the projection structure.
\textcite{fresca2022poddlrom} showed that a linear POD stage before the
autoencoder reduces training cost substantially without sacrificing accuracy.
For forecasting, \textcite{eivazi2021koopman} found a Koopman framework as
accurate as an LSTM in long-term statistics at lower cost, and
\textcite{morton2026scitech} found that transformers performed well on a
localized DLS latent space but poorly on a POD latent space of the same flow.
A common thread is that linear and nonlinear, global and local
representations are complementary. Hybrids that combine a linear or local
backbone with nonlinear corrections often perform best
\parencite{linot2020inertial,fresca2022poddlrom,diaz2024ddnm}.

\subsection{Baselines, benchmarks, and metrics}

\textcite{mcgreivy2024weak} documented that weak baselines and reporting biases
lead to systematic over-optimism in ML for fluid-related PDEs. For DR the
corresponding risks are comparing a tuned nonlinear model against an untuned
linear one, reporting training rather than held-out error, and choosing
latent sizes that favour one method. Based on the studies reviewed here, we
recommend the following practices.
\begin{enumerate}[leftmargin=2em,itemsep=2pt]
  \item Always report a POD baseline at the \emph{same} latent dimension and,
  for dynamics, a POD--based model with the same sequence architecture. Report
  POD error as a function of $r$, not at a single truncation.
  \item Evaluate reconstruction on held-out data that are separated in time
  from the training data and, where relevant, at unseen parameters or
  operating conditions (P1). Report training-data requirements (P4) alongside
  accuracy, as in \cref{fig:dls_compare}b.
  \item Report the full storage cost (latent coefficients \emph{and} basis or
  network parameters) when claiming compression, and the latent dimension seen
  by the dynamics model when claiming learnability.
  \item Beyond pointwise error, report second-order statistics (Reynolds
  stresses, turbulent kinetic energy), spectra, and, for forecasts, long-time
  statistics and error growth over rollouts. Small-scale fidelity is where
  methods differ most.
  \item Report variability over random initializations for learned
  representations.
\end{enumerate}
Common benchmark resources support such comparisons. They include the
Johns Hopkins turbulence databases \parencite{li2008jhtdb}, PDEBench
\parencite{takamoto2022pdebench}, The Well \parencite{ohana2024well}, and the
community challenge of \textcite{schmidt2026challenge}. The last has
separate tracks for compression, forecasting, and sensing of aerospace flows,
among them experimental time-resolved PIV of a Mach 0.6 cavity flow of the type
analysed by \textcite{zhang2020samm}.

\subsection{Practical guidance}

The choice of representation can be organized around four questions.
\begin{enumerate}[leftmargin=2em,itemsep=2pt]
  \item \emph{How fast does the POD spectrum decay?} If a few percent of the
  snapshots' rank captures the required energy on held-out data, POD (or DMD
  or SPOD for spectral content) is hard to beat and should be the default.
  \item \emph{Is the difficulty transport or multi-scale content?} Transport of a
  few coherent features favours transport-adapted linear methods or nonlinear
  autoencoders with small latent spaces. Broadband, multi-scale content over a
  large domain favours localized representations (domain decomposition, DLS,
  graph or local-token architectures) whose cost grows with the domain rather
  than with the number of distinct global patterns.
  \item \emph{How much data is there, and on what discretization?} Limited
  snapshots favour linear and local bases, since local bases draw many
  training samples from each snapshot. Unstructured or changing
  discretizations favour graph AEs and implicit neural representations.
  \item \emph{What will the latent space be used for?} Galerkin projection
  needs an explicit, differentiable decoder; control benefits from balanced or
  Koopman coordinates; long-horizon forecasting benefits from smooth,
  Markovian, and possibly localized coordinates; and generative modeling
  benefits from compact, well-regularized latent spaces.
\end{enumerate}
